\documentclass[10pt,a4paper]{amsart}
\usepackage[margin=19mm]{geometry}
\usepackage{amsmath,amssymb,mathtools}
\usepackage[scr=rsfs,cal=euler]{mathalfa}
\usepackage{xfrac}
\usepackage[unicode,hidelinks,bookmarks=false]{hyperref}
\newcommand{\ie}{i.e.\ }
\newcommand{\cf}{cf.\ }
\newcommand{\resp}{respectively\ }
\newcommand{\Subsection}{\subsection}
\providecommand{\nobreakdash}{\nobreak\hbox{-}}
\setlength{\emergencystretch}{2em}
\multlinegap0pt
\usepackage{bm}
\usepackage{enumerate}
\usepackage[shortlabels]{enumitem}
\usepackage{amssymb}
\usepackage{mathtools}
\usepackage{tikz}
\newtheorem{lemma}{Lemma}
\newcommand\bE{{\mathbb E}}
\newcommand\bN{{\mathbb N}}
\title{The temporal stochastic block model}
\author{Sofiya Burova, G\'abor Lugosi, Guillem Perarnau}
\date{Source: arXiv:2606.09488v1 (CC BY 4.0)}
\begin{document}
\maketitle

\noindent\emph{Source and licence.} The temporal stochastic block model. Version arXiv:2606.09488v1 (CC BY 4.0), distributed by arXiv under the Creative Commons Attribution 4.0 International licence, http://creativecommons.org/licenses/by/4.0/, as recorded in the arXiv licence metadata for this exact version and shown on the abstract page https://arxiv.org/abs/2606.09488v1. Full licence notice: \texttt{licenses/arxiv-2606.09488-CC-BY-4.0.txt}.\par\smallskip

\noindent\textit{Editorial scope.} Counted unit: the article's lemma lemma:nb\_paths (source0.tex lines 353--364). The article's complete original proof of it is delivered with it (source0.tex lines 366--392, one \texttt{proof} environment of 27 lines). No mathematics is written, repaired or re-derived: the statement and the proof are the article's own lines, in the article's own notation, and no retained line is substituted or re-flowed.  No further source-local context is needed: every cross-reference of every retained line resolves inside the two delivered spans.  Nothing was omitted from any retained span, so the package carries no omission record.  No retained line contains a citation, so the wrapper carries no bibliography and no bibliography entry.  No retained line contains a figure environment, an image-inclusion command or drawing code, so no image is delivered and the package declares no asset dependency.  Editorial formatting only, in addition: an AMS \texttt{amsart} preamble that restates the article's own declarations and macros used by the retained lines, namely the environments \texttt{lemma} on separate counters, together with the standard packages those lines need.  The retained environments print in this document as Lemma 1 in order of appearance, whereas the article prints the same items with the numbering of its own full document.  The complete native source of the article as distributed in the arXiv e-print for this exact version is bundled unchanged as \texttt{sources/202405/source0.tex}.
\begin{lemma}
\label{lemma:nb_paths}
For all $u,v \in V$ and for all $\ell \in \bN$,
\begin{equation*}
\begin{split}
\bE[X_\ell] &\leq  \frac{n (\log n)^\ell }{\ell!} \times (\boldsymbol{a}^T \cdot \Lambda^\ell \cdot \boldsymbol{1}), \\
\bE[Y_\ell(u)] &\leq \frac{(\log n)^\ell }{\ell!} \times (b_{\sigma(u)}^T \cdot \Lambda^\ell \cdot \boldsymbol{1}),\\
\bE[Z_\ell(u,v)] &\leq \frac{(\log n)^\ell }{a_{\sigma(v)}n\ell!} \times (b_{\sigma(u)}^T \cdot \Lambda^\ell \cdot b_{\sigma(v)}),
\end{split}
\end{equation*}
where $(b_i)_{i \in [k]}$ is the canonical basis of $\mathbb{R}^k$ and $\boldsymbol{1} = (1,\ldots,1)^T \in \mathbb{R}^k$. 
\end{lemma}

\begin{proof}
We first prove the third identity, since the other two can be easily derived from it. Fix $u,v\in V$ and denote by $Z^*_\ell(u,v)$ the number of paths (not necessarily increasing) of $\ell$ edges starting at $u$ and ending in $v$ in $G$. Notice that, since $G$ and $\pi$ are independent, for all $\ell \in \bN$,
$$\bE[Z_\ell(u,v)] = \frac{1}{\ell!} \bE[Z^*_\ell(u,v)]. $$
Hence, it suffices to show that
\begin{equation}
\label{ineq:induction}
\bE[Z^*_\ell(u,v) ]\leq \frac{(\log n)^\ell }{a_{\sigma(v)}n} \times (b_{\sigma(u)}^T \cdot \Lambda^\ell \cdot b_{\sigma(v)}). \end{equation}
We proceed by induction on $\ell \in \bN$. In the case $\ell =1$, the expected number of paths starting from $u$ and ending in $v$ is simply the probability of the edge $uv$ being present in the graph, that is,
$$\bE[Z^*_1(u,v)] = \frac{\lambda_{\sigma(u)\sigma(v)}\log n}{n} = \frac{\log n}{a_{\sigma(v)}n}\times (b^T_{\sigma(u)}\cdot \Lambda\cdot b_{\sigma(v)}). $$
Assuming that (\ref{ineq:induction}) holds for $\ell \in \bN$, we compute the expectation of $Z^*_{\ell+1}(u,v)$ by counting possible extensions of paths of length $\ell$ by an edge. For $z\in V\backslash\{u,v\}$, denote by $Z^*_\ell(u,z;v)$ the number of paths of length $\ell$ from $u$ to $z$, not containing $v$. Observe that $Z^*_{\ell}(u,z;v)$ and $Z^*_1(z,v)$ are independent. Thus,
\begin{equation*}
\begin{split}
    \bE[Z^*_{\ell+1}(u,v)] &= \sum_{z\in V\backslash\{u,v\}}\bE\left[  Z^*_\ell(u,z;v)  \right]\bE[Z^*_1(z,v)] \leq \sum_{z\in V}\bE\left[  Z^*_\ell(u,z)  \right]\bE[Z^*_1(z,v)]\\
    &\leq \sum_{z\in V} \frac{(\log n)^{\ell} }{a_{\sigma(z)}n} \times (b_{\sigma(u)}^T \cdot \Lambda^\ell \cdot b_{\sigma(z)}) \times \frac{ \Lambda_{\sigma(z)\sigma(v)}\log n}{a_{\sigma(v)}n}\\
    &= \frac{(\log n)^{\ell+1}}{a_{\sigma(v)}n} \sum_{i=1}^k \frac{|V_i|}{a_in} \times (\Lambda^\ell)_{\sigma(u)i} \Lambda_{i\sigma(v)} \\
    &= \frac{(\log n)^{\ell+1} }{a_{\sigma(v)}n} \times (b_{\sigma(u)}^T \cdot \Lambda^{\ell+1} \cdot b_{\sigma(v)}), 
\end{split}
\end{equation*}
where we used the induction hypothesis in the first inequality. This completes the proof of the upper bound on $\bE[Z_\ell(u,v)]$. Now, observe that
\begin{equation*}
    \begin{split}
        \bE[Y_\ell(u)] &\leq \sum_{v\in V} \bE[Z_\ell(u,v)] \leq \frac{(\log n)^\ell}{\ell!}\sum_{i=1}^k \frac{|V_i|}{a_i n} (\Lambda_{\sigma(u)i}) = \frac{(\log n)^\ell }{\ell!} \times (b_{\sigma(u)}^T \cdot \Lambda^\ell \cdot \boldsymbol{1}) , \\
        \bE[X_\ell] &\leq \sum_{u \in V}\bE[Y_{\ell}(u)] \leq \frac{(\log n)^\ell}{\ell!} \left(\sum_{i=1}^k |V_i|b_i \cdot \Lambda^\ell \cdot \boldsymbol{1} \right)  = \frac{n (\log n)^\ell }{\ell!} \times (\boldsymbol{a}^T \cdot \Lambda^\ell \cdot \boldsymbol{1}),
    \end{split}
\end{equation*}
which concludes the proof.
\end{proof}

\end{document}
